% Paragraph 2 (all storylines): the punishment dilemma.
% Round 2: ends with two routes out of the dilemma, which the storylines pick up.
The difficulty is sharpest for the norms that sustain cooperation under public information. Under these norms, a discriminator who refuses to help an individual in bad standing keeps its own good standing, because the refusal counts as justified punishment\cite{panchanathan2003,ohtsuki2004,okada2020dilemmas}. This exemption works only if the observer shares the donor's view of the recipient, which under private assessment is often not the case. An observer who still regards the recipient as good sees an unjustified refusal and condemns the discriminator. Disagreement about one individual thereby becomes disagreement about those who act towards it, and it spreads through the population\cite{hilbe2018,uchida2013,radzvilavicius2019}. Discriminators thus face a punishment dilemma. Withholding help from defectors protects cooperation against exploitation, yet each act of punishment risks their own reputation in the eyes of observers who disagree\cite{okada2020dilemmas,michelmata2024,pires2025}. The dilemma can be eased in two ways. Observers can be brought to agree about each other's standing, so that justified punishment is recognized as such, or reputations can be made robust enough that a single disputed judgement does not overturn them. In both cases, genuine defectors must still lose their standing.
